% =============================================================================
% CHAPTER 1 - INTRODUCTION
% =============================================================================

\chapter{INTRODUCTION}

\ifpdf
  \graphicspath{{introduction/figures/PNG}{introduction/figures/PDF}{introduction/figures}}
\else
  \graphicspath{{introduction/figures/EPS}{introduction/figures}}
\fi

\section{Background}

Despite their fluency, large language models (LLMs) are succeptible to hallucination, generating coherent and plausible text that is factually wrong~\citep{ji2023survey,huang2023survey}. Autoregressive language models generate text by sampling tokens from a learned distribution conditioned on prior context, with no mechanism for grounding outputs against external knowledge sources during decoding \citep{lewis2020rag}. Errors introduced at one decoding step can propagate to subsequent steps, and in knowledge-intensive tasks such errors manifest not as stylistic infelicities but as factual inaccuracies. This makes factuality a central limitation of LLMs in precisely the settings where correctness is most critical.

A knowledge graph (KG) stores facts as triplets $(\text{entity}, \text{relation}, \text{entity})$ and supports multi-hop reasoning over verifiable connections~\citep{bollacker2008freebase}. Benchmarks such as WebQSP~\citep{yih2016websqp} and ComplexWebQuestions~\citep{talmor2018cwq} test whether a system can trace these connections to answer natural language questions. Unconstrained LLMs, relying on parametric memory, frequently produce plausible but hallucinated reasoning trajectories over such graphs~\citep{he2021improving,lan2022survey}.

Retrieval-Augmented Generation (RAG) mitigates hallucination by injecting retrieved evidence into the context window \citep{lewis2020rag,izacard2021leveraging}. However, retrieved evidence functions as a soft prompt rather than a hard constraint on generation: models can disregard relevant retrieved context or generate claims unsupported by the retrieved material when internal parametric knowledge overrides external evidence \citep{asai2024selfrag}.

\newpage

\section{Problem Statement}

Constrained decoding takes a harder line by applying a validity mask over the output vocabulary. At each step, a constraint oracle forces the LLM to generate only tokens that correspond to valid paths in the KG~\citep{willard2023outlines,geng2023grammarconstrained}. Graph-Constrained Reasoning (GCR)~\citep{luo2025-graph-constrained-reasoning} achieves this with a pre-computed KG-Trie used as the oracle during beam search. The guarantee is structural in that, every generated token maps to a real graph edge. The problem is that the trie is built before decoding begins and never changes. It cannot adapt to the semantic intent of the question or the evolving reasoning chain, hence the LLM must still use its own judgment to filter out structurally valid but irrelevant paths~\citep{li2024-decoding-on-graphs,luo2025-graph-constrained-reasoning}.

Decoding on Graphs (DoG)~\citep{li2024-decoding-on-graphs} addresses this rigidity by expanding the trie dynamically during decoding. But DoG admits all topologically adjacent nodes without semantic filtering. The result is unnecessary computational overhead and continued vulnerability to off-target paths.

These gaps are the gaps that DCA-Trie is designed to close.By conditioning the valid token set on three inputs, the knowledge graph, the input question, and the partially generated reasoning chain, DCA-Trie combines structural constraints with dynamic semantic filtering, aiming to reduce constraint permissiveness without sacrificing the faithfulness guarantees that make constrained decoding valuable in the first place.

\section{Aim and Objectives}

The aim of this project is to design, implement, and evaluate a Dynamic Context-Aware Trie (DCA-Trie) framework for knowledge graph-constrained LLM generation. By conditioning each decoding step on the knowledge graph, the input question, and the partial reasoning chain to improve multi-hop KGQA accuracy while maintaining 100\% structural faithfulness to the underlying knowledge graph.

\newpage

The specific objectives of this project work are to:

\begin{enumerate}[label=\roman*.]
  \item Characterize the constraint permissiveness problem in existing constrained-decoding approaches to KG-grounded reasoning;
  \item Design a semantic relevance scoring mechanism and a constraint oracle for filtering structurally valid but semantically irrelevant paths;
  \item Implement a static filtering variant that applies the oracle at trie construction time, and evaluate its impact on false negative rate;
  \item Implement a dynamic variant in which the trie is updated during decoding, conditioned on the question and the partially generated reasoning chain; and
  \item Evaluate both variants against relevant baselines on standard knowledge-graph question answering benchmarks, and assess whether observed patterns generalize across benchmarks.
\end{enumerate}

\section{Tools and Facilities Used}

\subsection{Facilities}

The following facilities were utilized in the execution of this project:

\begin{enumerate}[label=\roman*.]
  \item NVIDIA GeForce RTX 4090 GPU (24\,GB VRAM) for all model inference and constrained decoding experiments.
  \item University library and online academic databases (ACM Digital Library, arXiv, Google Scholar, ACL Anthology) for literature review and citation verification.
  \item Shared team repository hosted on GitHub for version control, experiment logging, and reproducibility of all reported results.
\end{enumerate}

\subsection{Software Tools}

Table~\ref{tab:implemented-software-tools} presents the software tools used in this project.

\begin{table}[H]
  \centering
  \setlength{\tabcolsep}{4pt}
  \caption[Implemented Software Tools]{Implemented Software Tools}
  \label{tab:implemented-software-tools}
  \resizebox{\textwidth}{!}{%
    \begin{tabular}{|>{\raggedright\arraybackslash}p{0.22\textwidth}|>{\raggedright\arraybackslash}p{0.30\textwidth}|>{\raggedright\arraybackslash}p{0.32\textwidth}|}
      \hline
      \textbf{Tool}                     & \textbf{Use in Project}                                                                & \textbf{Rationale}                                                             \\
      \hline
      Python / HuggingFace Transformers & Model inference pipeline for Llama-3.1-8B, including the constrained decoding callback & Standard tokenization and generation utilities for the released GCR checkpoint \\
      \hline
      GCR Framework                     & Primary baseline; KG-Trie construction and graph-constrained decoding                  & Provides the trie-based decoding mechanism DCA-Trie extends                    \\
      \hline
      PyTorch                           & Deep learning framework underpinning inference and decoding logic                      & Supports SDPA attention and bfloat16 inference                                 \\
      \hline
      marisa\_trie                      & KG-Trie construction and prefix-based token lookup                                     & $O(1)$ prefix lookup for the valid-token mask at each decoding step            \\
      \hline
      NetworkX                          & Graph construction and path enumeration from question entities                         & Builds directed graphs from Freebase triplets                                  \\
      \hline
      Freebase / WebQSP / CWQ           & Knowledge graph source and KGQA benchmarks                                             & Standard benchmarks for comparison with prior work                             \\
      \hline
      NVIDIA RTX 4090 (24GB)            & GPU for all inference experiments                                                      & Sufficient VRAM for bfloat16 inference of an 8B-parameter model                \\
      \hline
      VS Code, Git / GitHub             & Development environment and version control                                            & Coding, version tracking, and reproducibility                                  \\
      \hline
    \end{tabular}
  }
\end{table}

\section{Scope and Delimitation}

This project targets the permissiveness problem in GCR and DoG: existing constraint oracles admit every structurally reachable path regardless of relevance. Evaluation is restricted to multi-hop KGQA over Freebase benchmarks (WebQSP and CWQ). All experiments use GCR's publicly released Llama-3.1-8B checkpoint~\citep{luo2025-graph-constrained-reasoning}; no fine-tuning is performed. Generalization to other KGs or non-KGQA tasks is identified as future work. All code and evaluation datasets will be made publicly available.

\section{Organization of the Project Report}

The rest of the report is organised as follows. Chapter 2 reviews literature on hallucination, knowledge graph question answering, and constrained decoding. Chapter 3 presents the research methodology, formalising DCA-Trie and its two variants. Chapter 4 reports results on WebQSP and CWQ, including ablation and error analysis. Chapter 5 presents conclusions and recommendations for future work.